\documentclass[../../main.tex]{subfiles}
\graphicspath{{\subfix{../../image/}}}

\begin{document}

\section{随机变量的函数及其分布}\label{section:ch03-随机变量的函数及其分布}

\subsection*{一、Borel 函数与随机变量的函数}

在许多问题中需要计算随机变量的函数的分布律，例如在统计物理中，已知分子运动速度 $\xi$ 的分布，要求其动能 $\eta=\dfrac{1}{2}m\xi^{2}$ 的分布律。这类问题既普遍而又重要，接下来我们就要讨论它。

这类问题较为一般的提法是：若 $\xi$ 是随机变量，求 $\eta=g(\xi)$ 的分布律。

为了使 $\eta$ 有分布律可言，当然要求 $\eta$ 是随机变量，因此对函数 $y=g(x)$ 也必须有一定的要求，为此我们提出如下定义。

\begin{definition}[Borel 可测函数]\label{definition:ch03-Borel可测函数}
设 $y=g(x)$ 是 $\mathbf{R}^{1}$ 到 $\mathbf{R}^{1}$ 上的一个映照，若对于一切 $\mathbf{R}^{1}$ 中的 Borel 点集 $B_1$ 均有
\begin{align}
\{x:g(x)\in B_1\}\in\mathscr{B}_1,\label{align:ch03-Borel可测}
\end{align}
其中 $\mathscr{B}_1$ 为 $\mathbf{R}^{1}$ 上 Borel $\sigma$ 域，则称 $g(x)$ 是一元 Borel（可测）函数。
\end{definition}

Borel 函数是很广泛的一类函数，我们所碰到的大部分函数都是 Borel 函数，特别地，已知所有的连续函数及单调函数都是 Borel 函数。

若 $\xi$ 是概率空间 $(\Omega,\mathscr{F},P)$ 上的随机变量，而 $g(x)$ 是一元 Borel 函数，则 $g(\xi)$ 是 $(\Omega,\mathscr{F},P)$ 上的随机变量。事实上，对一切 $B_1\in\mathscr{B}_1$ 有
\begin{align}
\{\omega:g(\xi(\omega))\in B_1\}=\{\omega:\xi(\omega)\in g^{-1}(B_1)\}\in\mathscr{F}.\label{align:ch03-函数随机变量}
\end{align}
这里 $g^{-1}(B_1)=\{x:g(x)\in B_1\}$，由\eqref{align:ch03-Borel可测}知它是一维 Borel 点集，再根据随机变量定义中的\eqref{align:ch03-随机变量可测}式即得\eqref{align:ch03-函数随机变量}。

有时，我们还要考虑随机向量的函数，例如 $\eta=\xi_1+\xi_2$ 就是随机向量 $(\xi_1,\xi_2)$ 的函数。一般地，要研究 $\eta=g(\xi_1,\cdots,\xi_n)$ 及其分布律。这时对 $n$ 元函数 $y=g(x_1,\cdots,x_n)$ 也要有相应的要求。

\begin{definition}[多元 Borel 可测函数]\label{definition:ch03-多元Borel可测函数}
设 $y=g(x_1,\cdots,x_n)$ 是 $\mathbf{R}^{n}$ 到 $\mathbf{R}^{1}$ 上的一个映照，若对一切 $\mathbf{R}^{1}$ 中的 Borel 点集 $B_1$ 均有
\begin{align}
\{(x_1,\cdots,x_n):g(x_1,\cdots,x_n)\in B_1\}\in\mathscr{B}_n,\label{align:ch03-多元Borel可测}
\end{align}
其中 $\mathscr{B}_n$ 为 $\mathbf{R}^{n}$ 上 Borel $\sigma$ 域，则称 $g(x_1,\cdots,x_n)$ 为 $n$ 元 Borel（可测）函数。
\end{definition}

同样地，若 $(\xi_1,\cdots,\xi_n)$ 是 $(\Omega,\mathscr{F},P)$ 上的随机向量，而 $g(x_1,\cdots,x_n)$ 是 $n$ 元 Borel 函数，则 $g(\xi_1,\cdots,\xi_n)$ 是 $(\Omega,\mathscr{F},P)$ 上的随机变量。

这个事实证明如下：
\begin{align*}
\{\omega:g(\xi_1(\omega),\cdots,\xi_n(\omega))\in B_1\}
=\{\omega:(\xi_1(\omega),\cdots,\xi_n(\omega))\in g^{-1}(B_1)\}\in\mathscr{F}.
\end{align*}
其中 $g^{-1}(B_1)$ 是\eqref{align:ch03-多元Borel可测}中点集，它是 $n$ 维 Borel 点集，而按\eqref{align:ch03-随机向量Borel可测}即得最后的关系式。

更一般地，还可以研究 $n$ 维随机向量 $(\xi_1,\cdots,\xi_n)$ 的 $m$ 个函数 $g_1(\xi_1,\cdots,\xi_n),\cdots,g_m(\xi_1,\cdots,\xi_n)$，这里 $g_1,\cdots,g_m$ 都是 $n$ 元 Borel 函数，这是一个 $m$ 维随机向量，需要求出它们的分布。

一般地来说，对于离散型随机变量，求它的函数的分布并不很困难。例如：若 $\xi$ 的分布列为
\begin{align*}\begin{array}{c|cccccc}
\xi & x_1 & x_2 & \cdots & x_n & \cdots\\
\hline
P & p_1 & p_2 & \cdots & p_n & \cdots
\end{array}\end{align*}
则 $g(\xi)$ 的分布列可由下法得到，列出
\begin{align}
\begin{array}{c|cccccc}
g(\xi) & g(x_1) & g(x_2) & \cdots & g(x_n) & \cdots\\
\hline
P & p_1 & p_2 & \cdots & p_n & \cdots
\end{array}\label{align:ch03-函数分布列}
\end{align}
当然这里可能有某些 $g(x_i)$ 相等，把它们作适当并项即可。

\paragraph{[离散卷积公式]}若 $\xi$ 与 $\eta$ 是相互独立的随机变量，它们都取非负整数值，其概率分布分别为 $\{a_k\}$ 及 $\{b_k\}$，下面我们来计算随机变量 $\zeta=\xi+\eta$ 的概率分布。因为
\begin{align*}
\{\zeta=r\}=\{\xi=0,\eta=r\}+\{\xi=1,\eta=r-1\}+\cdots+\{\xi=r,\eta=0\}.
\end{align*}
利用独立性的假定得到
\begin{align}
c_r=P\{\zeta=r\}=a_0b_r+a_1b_{r-1}+\cdots+a_rb_0,\qquad r=0,1,2,\cdots.\label{align:ch03-离散卷积}
\end{align}
这就是求独立随机变量和的分布的公式——\textbf{离散卷积公式}。

下面的讨论主要对连续型随机变量进行。我们将由简到繁逐步深入。

\subsection*{二、单个随机变量的函数的分布律}

这里的一般问题是：已知随机变量 $\xi$ 的分布函数 $F(x)$ 或密度函数 $p(x)$，要求 $\eta=g(\xi)$ 的分布函数 $G(y)$ 或密度函数 $q(y)$。

由\eqref{align:ch03-密度Borel概率}可知
\begin{align}
G(y)=P\{\eta<y\}=P\{g(\xi)<y\}=\int_{\{x:g(x)<y\}}p(x)\,\mathrm{d}x.\label{align:ch03-函数分布积分}
\end{align}
上述积分计算的难易既与被积函数即 $\xi$ 的密度函数 $p(x)$ 的表达式有关，更与积分区域 $\{x:g(x)<y\}$ 的形状相关，差别很大，因此这类问题通常采用个案处理的方式，但在方法上大体可分为直接法与变换法两类。

直接法通过把 $\{g(\xi)<y\}$ 直接化为关于 $\xi$ 的等价事件而求得 $\eta$ 的分布函数或密度函数。

当 $g(x)$ 为单调函数时，问题相当简单。例如当 $g(x)$ 严格单调上升时，
\begin{align*}G(y)=P\{\eta<y\}=P\{g(\xi)<y\}=P\{\xi<g^{-1}(y)\}=F(g^{-1}(y)).\end{align*}

\begin{example}[线性变换]\label{example:ch03-线性变换}
若随机变量 $\xi$ 有密度函数 $p(x)$，而 $\eta=a\xi+b$，这里 $a\neq 0$。求 $\eta$ 的密度函数 $q(y)$。

\paragraph*{解}分别记 $\xi$ 与 $\eta$ 的分布函数为 $F(x)$ 及 $G(y)$，显然有
\begin{align*}
G(y)=P\{\eta<y\}=P\{a\xi+b<y\}=
\begin{cases}
P\left\{\xi<\dfrac{y-b}{a}\right\}=F\left(\dfrac{y-b}{a}\right), & \text{若 }a>0,\\[4pt]
P\left\{\xi>\dfrac{y-b}{a}\right\}=1-F\left(\dfrac{y-b}{a}\right), & \text{若 }a<0.
\end{cases}
\end{align*}
因此
\begin{align*}q(y)=
\begin{cases}
\dfrac{1}{a}p\left(\dfrac{y-b}{a}\right), & \text{若 }a>0,\\[4pt]
-\dfrac{1}{a}p\left(\dfrac{y-b}{a}\right), & \text{若 }a<0,
\end{cases}\end{align*}
或统一起来，写成
\begin{align}
q(y)=\frac{1}{|a|}p\left(\frac{y-b}{a}\right).\label{align:ch03-线性变换密度}
\end{align}

把例\ref{example:ch03-线性变换}的结果应用到正态分布的场合。若 $\xi$ 服从 $N(\mu,\sigma^{2})$，则 $\zeta=\dfrac{\xi-\mu}{\sigma}$ 的密度函数
\begin{align*}q(z)=\sigma\cdot\frac{1}{\sqrt{2\pi}\sigma}\exp\left\{-\frac{\left[\sigma\left(z+\frac{\mu}{\sigma}\right)-\mu\right]^{2}}{2\sigma^{2}}\right\}=\frac{1}{\sqrt{2\pi}}\mathrm{e}^{-z^{2}/2},\end{align*}
即 $\zeta$ 服从 $N(0,1)$，这是我们早已提到过的。
\end{example}

\begin{example}[对数正态分布]\label{example:ch03-对数正态}
若 $\xi\sim N(\mu,\sigma^{2})$，求 $\eta=\mathrm{e}^{\xi}$ 的密度函数。

\paragraph*{解}当 $y>0$ 时，
\begin{align*}
P\{\eta<y\}=P\{\mathrm{e}^{\xi}<y\}=P\{\xi<\ln y\}=\int_{-\infty}^{\ln y}\frac{1}{\sqrt{2\pi}\sigma}\mathrm{e}^{-\frac{(x-\mu)^{2}}{2\sigma^{2}}}\,\mathrm{d}x.
\end{align*}
所以，$\eta$ 的密度函数为
\begin{align}
q(y)=\frac{1}{\sqrt{2\pi}\sigma y}\mathrm{e}^{-\frac{(\ln y-\mu)^{2}}{2\sigma^{2}}},\qquad y>0.\label{align:ch03-对数正态密度}
\end{align}

$\eta$ 的对数即 $\ln\eta=\xi$ 服从正态分布，故称 $\eta$ 所服从的分布为\textbf{对数正态分布}。

对数正态变量取非负值，又能通过正态分布进行概率计算，很适合做某些现象的数学模型。当代金融学用对数正态分布取代正态分布作为资产价格分布建立起了十分漂亮而合理的理论。另外，销售量，元件寿命等也已普遍使用对数正态分布作为模型。这个分布的重要性正在提高。

显然，上述做法可推广到 $g(x)$ 分段单调的场合。
\end{example}

\begin{example}[平方变换]\label{example:ch03-平方变换}
若 $\zeta\sim N(0,1)$，求 $\eta=\zeta^{2}$ 的密度函数。

\paragraph*{解}当 $y\leqslant 0$ 时，$G(y)=P\{\eta<y\}=0$，显然，此时 $q(y)=0$。当 $y>0$，
\begin{align*}
G(y)=P\{\eta<y\}=P\{\zeta^{2}<y\}=P\{-\sqrt{y}<\zeta<\sqrt{y}\}
=\int_{-\sqrt{y}}^{\sqrt{y}}\frac{1}{\sqrt{2\pi}}\mathrm{e}^{-\frac{x^{2}}{2}}\,\mathrm{d}x
=2\int_{0}^{\sqrt{y}}\frac{1}{\sqrt{2\pi}}\mathrm{e}^{-\frac{x^{2}}{2}}\,\mathrm{d}x.
\end{align*}
因此 $\eta=\zeta^{2}$ 的密度函数为
\begin{align}
q(y)=\frac{1}{\sqrt{2\pi}}y^{-\frac{1}{2}}\mathrm{e}^{-\frac{y}{2}},\qquad y>0.\label{align:ch03-卡方1密度}
\end{align}
顺便指出，\eqref{align:ch03-卡方1密度}是下列分布当 $n=1$ 时的特例。

\paragraph{[$\chi^{2}$ 分布]}具有密度函数
\begin{align}
p(x)=\frac{1}{2^{n/2}\Gamma\!\left(\frac{n}{2}\right)}x^{\frac{n}{2}-1}\mathrm{e}^{-\frac{x}{2}},\qquad x>0\label{align:ch03-卡方密度}
\end{align}
的分布称为具有自由度 $n$ 的 $\chi^{2}$ 分布。

$\chi^{2}$ 分布在数理统计中有重要应用。与\eqref{align:ch03-gamma密度}中的 $\Gamma$ 分布比较，我们就可以知道 $\chi^{2}$ 分布是它的特例，其中取 $r=\dfrac{n}{2}$，$\lambda=\dfrac{1}{2}$，即为 $\Gamma\!\left(\dfrac{n}{2},\dfrac{1}{2}\right)$。以后简记 $\xi$ 服从自由度 $n$ 的 $\chi^{2}$ 分布为 $\xi\sim\chi_n^{2}$。
\end{example}

下例处理更为一般的场合。

\begin{example}[正弦变换]\label{example:ch03-正弦变换}
若 $\xi$ 为连续型随机变量，其分布密度为 $p(x)$，求 $\eta=\sin\xi$ 的分布。

\paragraph*{解}记 $G(y)=P\{\sin\xi<y\}$。当 $y\leqslant -1$ 时，$G(y)=0$，当 $y>1$ 时，$G(y)=1$。当 $-1<y\leqslant 1$ 时（图\ref{figure:ch03-正弦分布}），
\begin{align*}
G(y)=P\{\sin\xi<y\}=\sum_{k=-\infty}^{\infty}P\{2k\pi-(\pi+\sin^{-1}y)\leqslant\xi<2k\pi+\sin^{-1}y\}
=\int_{\{x:\sin x<y\}}p(x)\,\mathrm{d}x.
\end{align*}
\end{example}

\begin{figure}[H]
\centering
\includegraphics[width=0.6\textwidth]{ch03-fig-sin.png}
\caption{$\eta=\sin\xi$ 的分布}
\label{figure:ch03-正弦分布}
\end{figure}

变换法利用积分的变数代换，对密度函数得到一般的计算公式。这时对函数 $g(x)$ 有较强的要求，但其结果可推广到多维场合，因此是较为有力的方法。下面仍依单调与分段单调两种情况进行讨论。

若 $\xi$ 为连续型随机变量，其密度函数为 $p(x)$，而 $\eta=g(\xi)$，对其密度函数 $q(y)$ 有如下结果：

(1) 若 $g(x)$ 严格单调，其反函数 $g^{-1}(y)$ 有连续导函数，则 $\eta=g(\xi)$ 是具有密度函数\footnote{这里和下面都约定，对使反函数无意义的 $y$，密度函数定义为 $0$。}
\begin{align}
p[g^{-1}(y)]\,|[g^{-1}(y)]'|\label{align:ch03-单调变换密度}
\end{align}
的连续型随机变量。

\begin{proof}
对于任一实数 $a$，记使 $g(x)<a$ 成立的 $x$ 的值的范围为 $E(a)$，则
\begin{align*}
P\{\eta<a\}=P\{g(\xi)<a\}=P\{\xi\in E(a)\}
=\int_{E(a)}p(x)\,\mathrm{d}x=\int_{-\infty}^{a}p[g^{-1}(y)]\,|[g^{-1}(y)]'|\,\mathrm{d}y.
\end{align*}
这里用了积分的变数代换。由上式知道 $\eta$ 为连续型随机变量，其密度函数为\eqref{align:ch03-单调变换密度}。
\end{proof}

显然可利用这个结果直接计算例\ref{example:ch03-线性变换}及例\ref{example:ch03-对数正态}中的密度函数。下面是另一个例子。

\begin{example}[Cauchy 分布]\label{example:ch03-cauchy}
若 $\theta$ 服从 $\left[-\dfrac{\pi}{2},\dfrac{\pi}{2}\right]$ 的均匀分布，$\psi=\operatorname{tg}\theta$，试求 $\psi$ 的密度函数 $q(y)$。

\paragraph*{解}记 $y=\operatorname{tg}x$，则 $x=\operatorname{tg}^{-1}y$，$\dfrac{\mathrm{d}x}{\mathrm{d}y}=\dfrac{1}{1+y^{2}}$，因此由\eqref{align:ch03-单调变换密度}知
\begin{align}
q(y)=\frac{1}{\pi}\cdot\frac{1}{1+y^{2}},\qquad -\infty<y<\infty.\label{align:ch03-cauchy密度}
\end{align}
由\eqref{align:ch03-cauchy密度}定义的分布称为 Cauchy 分布，它是概率论中有名的分布之一。
\end{example}

(2) 若 $g(x)$ 在不相重叠的区间 $I_1,I_2,\cdots$ 上逐段严格单调，其反函数分别为 $h_1(y),h_2(y),\cdots$ 而且 $h_1'(y),h_2'(y),\cdots$ 均为连续函数，那么 $\eta=g(\xi)$ 是连续型随机变量，其密度函数为
\begin{align}
p[h_1(y)]\,|h_1'(y)|+p[h_2(y)]\,|h_2'(y)|+\cdots.\label{align:ch03-分段单调密度}
\end{align}

\begin{proof}
给定实数 $a$，以 $E_i(a)$ 记 $I_i$ 中使 $g(x)<a$ 成立的 $x$ 的集合，显然诸 $E_i(a)$ 不相交，而
\begin{align*}
P\{\eta<a\}=P\{g(\xi)<a\}=P\left\{\xi\in\sum_{i}E_i(a)\right\}
=\sum_{i}\int_{E_i(a)}p(x)\,\mathrm{d}x
=\sum_{i}\int_{-\infty}^{a}p[h_i(y)]\,|h_i'(y)|\,\mathrm{d}y
=\int_{-\infty}^{a}\sum_{i}p[h_i(y)]\,|h_i'(y)|\,\mathrm{d}y.
\end{align*}
因此 $\eta$ 是具有密度函数\eqref{align:ch03-分段单调密度}的连续型随机变量。
\end{proof}

这结果包含了单调的场合\eqref{align:ch03-单调变换密度}，此外，它也可以用于例\ref{example:ch03-平方变换}及例\ref{example:ch03-正弦变换}，并得出相同答案。

这里，我们要再提到一个有趣而重要的事实。

\paragraph{[均匀分布的特殊地位]}若随机变量 $\xi$ 的分布函数为 $F(x)$，因为 $F(x)$ 是非降函数，对任意 $0\leqslant y\leqslant 1$，可定义
\begin{align}
F^{-1}(y)=\inf\{x:F(x)>y\}\label{align:ch03-逆分布函数}
\end{align}
作为 $F(x)$ 的反函数。

下面考察随机变量 $\theta=F(\xi)$ 的分布，这里 $F(x)$ 是连续函数。对 $0\leqslant x\leqslant 1$，
\begin{align}
P\{\theta<x\}=P\{F(\xi)<x\}=P\{\xi<F^{-1}(x)\}=F(F^{-1}(x))=x.\label{align:ch03-概率积分变换}
\end{align}
即 $\theta=F(\xi)$ 服从 $[0,1]$ 均匀分布。这个结论在统计中起重要作用。

反之，若 $\theta$ 服从 $[0,1]$ 均匀分布，对任意分布函数 $F(x)$，令
\begin{align}
\xi=F^{-1}(\theta)\label{align:ch03-逆变换采样}
\end{align}
则
\begin{align*}P\{\xi<x\}=P\{F^{-1}(\theta)<x\}=P\{\theta<F(x)\}=F(x).\end{align*}
因此 $\xi$ 是服从分布函数 $F(x)$ 的随机变量。

这样，只要我们能产生 $[0,1]$ 中均匀分布的随机变量的样本（观察值），那么我们也就能通过\eqref{align:ch03-逆变换采样}产生分布函数为 $F(x)$ 的随机变量的样本，这结论在 Monte Carlo 方法中具有基本的重要性。通常的做法是利用数学或物理的方法产生 $[0,1]$ 中均匀分布随机变量的样本（称为\textbf{均匀分布随机数}），再利用变换\eqref{align:ch03-逆变换采样}得到任意分布 $F(x)$ 的随机数。

\paragraph*{随机变量的存在性定理}利用上述结果，我们可以给下面定理一个构造性的证明。

\begin{theorem}[随机变量的存在性]\label{theorem:ch03-随机变量存在性}
若 $F(x)$ 是左连续的单调不减函数，且 $F(-\infty)=0,F(+\infty)=1$，则存在一个概率空间 $(\Omega,\mathscr{F},P)$ 及其上的随机变量 $\xi(\omega)$，使 $\xi(\omega)$ 的分布函数正好是 $F(x)$。
\end{theorem}

\begin{proof}
取 $\Omega=[0,1]$，再取 $\mathscr{F}$ 为 $[0,1]$ 中 Borel 点集全体，而 $P$ 取为直线上的 Lebesgue 测度（它是长度概念的推广，但对一切 Borel 点集都有定义）。定义 $\theta(\omega)=\omega$，则 $\theta(\omega)$ 是 $(\Omega,\mathscr{F},P)$ 上的随机变量，又对一切 $0\leqslant x\leqslant 1$，
\begin{align*}P\{\theta(\omega)<x\}=P\{\omega\in[0,x)\}=x.\end{align*}
因此 $\theta(\omega)$ 服从 $[0,1]$ 上均匀分布。

再利用\eqref{align:ch03-逆分布函数}定义 $F^{-1}(y)$，当然它也是单调函数，从而是 Borel 函数，令
\begin{align*}\xi(\omega)=F^{-1}(\theta(\omega)).\end{align*}
则 $\xi(\omega)$ 是 $(\Omega,\mathscr{F},P)$ 上的随机变量，而且仿上段讨论可知，它的分布函数正好是 $F(x)$。
\end{proof}

\subsection*{三、随机向量的函数的分布律}

若 $\eta=g(\xi_1,\cdots,\xi_n)$，而 $(\xi_1,\cdots,\xi_n)$ 的密度函数为 $p(x_1,\cdots,x_n)$，则同上面一样讨论可以得到
\begin{align}
G(y)=P\{\eta<y\}=\int_{\{g(x_1,\cdots,x_n)<y\}}\cdots\int p(x_1,\cdots,x_n)\,\mathrm{d}x_1\cdots\mathrm{d}x_n.\label{align:ch03-向量函数分布}
\end{align}

我们看一些例子。

\paragraph{[和的分布]}若 $\eta=\xi_1+\xi_2$，而 $(\xi_1,\xi_2)$ 的密度函数为 $p(x_1,x_2)$，则
\begin{align}
\begin{aligned}
G(y)=P\{\eta<y\}&=\iint_{\{x_1+x_2<y\}}p(x_1,x_2)\,\mathrm{d}x_1\,\mathrm{d}x_2\\
&=\int_{-\infty}^{\infty}\int_{-\infty}^{y-x_1}p(x_1,x_2)\,\mathrm{d}x_1\,\mathrm{d}x_2.
\end{aligned}\label{align:ch03-和的分布}
\end{align}
特别当 $\xi_1,\xi_2$ 相互独立时，有 $p(x_1,x_2)=p_1(x_1)p_2(x_2)$，这里 $p_1(x_1)$ 为 $\xi_1$ 的密度函数，$p_2(x_2)$ 为 $\xi_2$ 的密度函数。代入\eqref{align:ch03-和的分布}得
\begin{align*}
G(y)&=\int_{-\infty}^{\infty}\left[\int_{-\infty}^{y-x_1}p_1(x_1)p_2(x_2)\,\mathrm{d}x_2\right]\mathrm{d}x_1\\
&=\int_{-\infty}^{\infty}\left[\int_{-\infty}^{y}p_1(x_1)p_2(z-x_1)\,\mathrm{d}z\right]\mathrm{d}x_1\\
&=\int_{-\infty}^{y}\left[\int_{-\infty}^{\infty}p_1(x_1)p_2(z-x_1)\,\mathrm{d}x_1\right]\mathrm{d}z.
\end{align*}
因此 $\eta$ 的密度函数为
\begin{align}
q(y)=\int_{-\infty}^{\infty}p_1(u)p_2(y-u)\,\mathrm{d}u.\label{align:ch03-连续卷积1}
\end{align}
也可写为
\begin{align}
q(y)=\int_{-\infty}^{\infty}p_1(y-u)p_2(u)\,\mathrm{d}u.\label{align:ch03-连续卷积2}
\end{align}
\eqref{align:ch03-连续卷积1}或\eqref{align:ch03-连续卷积2}称为\textbf{卷积公式}，在概率论中相当重要。

\paragraph{[商的分布]}若 $\eta=\dfrac{\xi_1}{\xi_2}$，而 $(\xi_1,\xi_2)$ 的密度函数为 $p(x_1,x_2)$，则
\begin{align}
\begin{aligned}
G(x)=P\{\eta<x\}=P\left\{\frac{\xi_1}{\xi_2}<x\right\}
&=\iint_{\{x_1/x_2<x\}}p(x_1,x_2)\,\mathrm{d}x_1\,\mathrm{d}x_2\\
&=\int_{0}^{\infty}\left[\int_{-\infty}^{zx}p(y,z)\,\mathrm{d}y\right]\mathrm{d}z+\int_{-\infty}^{0}\left[\int_{zx}^{\infty}p(y,z)\,\mathrm{d}y\right]\mathrm{d}z.
\end{aligned}\label{align:ch03-商的分布}
\end{align}
$\eta$ 的密度函数为
\begin{align}
q(x)=\int_{0}^{\infty}p(zx,z)z\,\mathrm{d}z-\int_{-\infty}^{0}p(zx,z)z\,\mathrm{d}z
=\int_{-\infty}^{\infty}|z|\,p(zx,z)\,\mathrm{d}z.\label{align:ch03-商的密度}
\end{align}

\paragraph*{[关于顺序统计量的若干分布]}若 $\xi_1,\xi_2,\cdots,\xi_n$ 是相互独立的随机变量，具有相同的分布函数 $F(x)$ 和密度函数 $p(x)$，而 $\xi_n^{*}$ 及 $\xi_1^{*}$ 相当于把 $\xi_1,\xi_2,\cdots,\xi_n$ 按大小顺序重新排列为
\begin{align}
\xi_1^{*}\leqslant\xi_2^{*}\leqslant\cdots\leqslant\xi_n^{*}\label{align:ch03-顺序统计量}
\end{align}
的末项及首项，它们在统计中有重要应用。下面讨论几种与它们有关的分布。

首先求极大值 $\xi_n^{*}$ 的分布函数，
\begin{align}
\begin{aligned}
P\{\xi_n^{*}<x\}&=P\{\max(\xi_1,\xi_2,\cdots,\xi_n)<x\}\\
&=P\{\xi_1<x,\xi_2<x,\cdots,\xi_n<x\}\\
&=P\{\xi_1<x\}\cdot P\{\xi_2<x\}\cdots P\{\xi_n<x\}\\
&=[F(x)]^{n}.
\end{aligned}\label{align:ch03-极大值分布}
\end{align}
其次求极小值 $\xi_1^{*}$ 的分布函数，注意到
\begin{align*}
P\{\xi_1^{*}\geqslant x\}=P\{\min(\xi_1,\xi_2,\cdots,\xi_n)\geqslant x\}
=P\{\xi_1\geqslant x,\xi_2\geqslant x,\cdots,\xi_n\geqslant x\}
=P\{\xi_1\geqslant x\}P\{\xi_2\geqslant x\}\cdots P\{\xi_n\geqslant x\}
=[1-F(x)]^{n}.
\end{align*}
因此
\begin{align}
P\{\xi_1^{*}<x\}=1-[1-F(x)]^{n}.\label{align:ch03-极小值分布}
\end{align}
值得提醒，\eqref{align:ch03-极小值分布}式的导出，是公式\eqref{align:ch03-和的分布}的生动应用。

进一步，讨论 $(\xi_1^{*},\xi_n^{*})$ 的联合分布。

记 $G(x,y)=P\{\xi_1^{*}<x,\xi_n^{*}<y\}$。

若 $x\geqslant y$，则
\begin{align}
G(x,y)=P\{\xi_1^{*}<x,\xi_n^{*}<y\}=P\{\xi_n^{*}<y\}=[F(y)]^{n}.\label{align:ch03-极值联合xgey}
\end{align}
若 $x<y$，则
\begin{align}
\begin{aligned}
G(x,y)&=P\{\xi_1^{*}<x,\xi_n^{*}<y\}\\
&=P\{\xi_n^{*}<y\}-P\{\xi_1^{*}\geqslant x,\xi_n^{*}<y\}\\
&=[F(y)]^{n}-[F(y)-F(x)]^{n}.
\end{aligned}\label{align:ch03-极值联合xlty}
\end{align}
其联合密度函数为
\begin{align}
q(x,y)=
\begin{cases}
0, & x\geqslant y,\\
n(n-1)[F(y)-F(x)]^{n-2}p(x)p(y), & x<y.
\end{cases}\label{align:ch03-极值联合密度}
\end{align}

最后，我们来求极差 $R=\xi_n^{*}-\xi_1^{*}$ 的分布密度函数 $f_R(r)$，显然对 $r\leqslant 0,f_R(r)=0$，若 $r>0$，则
\begin{align*}
P\{R<r\}=\iint_{\{y-x<r\}}q(x,y)\,\mathrm{d}x\,\mathrm{d}y
=\int_{-\infty}^{\infty}\left[\int_{-\infty}^{x+r}q(x,y)\,\mathrm{d}y\right]\mathrm{d}x
=\int_{-\infty}^{\infty}\left[\int_{-\infty}^{r}q(x,x+z)\,\mathrm{d}z\right]\mathrm{d}x
=\int_{-\infty}^{r}\left[\int_{-\infty}^{\infty}q(x,x+z)\,\mathrm{d}x\right]\mathrm{d}z.
\end{align*}
因此
\begin{align}
f_R(r)=\int_{-\infty}^{\infty}q(x,x+r)\,\mathrm{d}x
=n(n-1)\int_{-\infty}^{\infty}[F(x+r)-F(x)]^{n-2}p(x)p(x+r)\,\mathrm{d}x.\label{align:ch03-极差密度}
\end{align}

极值分布在统计中常被用到。在实际应用中，极值分布与“百年一遇”等概念经常出现在灾害性天气预报中，例如暴雨，洪水预报，以及水库、桥梁等大型工程建筑规范中。

\subsection*{四、随机向量的变换}

若 $(\xi_1,\cdots,\xi_n)$ 的密度函数为 $p(x_1,\cdots,x_n)$，求 $\eta_1=g_1(\xi_1,\cdots,\xi_n),\cdots,\eta_m=g_m(\xi_1,\cdots,\xi_n)$ 的分布。这时有
\begin{align}
G(y_1,\cdots,y_m)=P\{\eta_1<y_1,\cdots,\eta_m<y_m\}
=\int_{\substack{g_1(x_1,\cdots,x_n)<y_1\\\cdots\\g_m(x_1,\cdots,x_n)<y_m}}\cdots\int p(x_1,\cdots,x_n)\,\mathrm{d}x_1\cdots\mathrm{d}x_n.\label{align:ch03-向量变换分布}
\end{align}

显然，这是最一般的场合。当 $m=1$ 时便是随机向量的函数的情形，当 $m=n=1$ 时得到单个随机变量的函数的情形。下面考虑另一个重要的特殊情形，即当 $(\xi_1,\cdots,\xi_n)$ 与 $(\eta_1,\cdots,\eta_m)$ 有一一对应变换关系时，当然这时 $n=m$ 必须成立。

如果对 $y_i=g_i(x_1,\cdots,x_n)$，$i=1,2,\cdots,n$，存在唯一的反函数 $x_i(y_1,\cdots,y_n)=x_i$（$i=1,\cdots,n$），而且 $(\eta_1,\cdots,\eta_n)$ 的密度函数为 $q(y_1,\cdots,y_n)$，那么
\begin{align}
G(y_1,\cdots,y_n)=\int_{\substack{u_1<y_1\\\cdots\\u_n<y_n}}\cdots\int q(u_1,\cdots,u_n)\,\mathrm{d}u_1\cdots\mathrm{d}u_n.\label{align:ch03-逆变换分布}
\end{align}
比较 $m=n$ 时的\eqref{align:ch03-向量变换分布}与\eqref{align:ch03-逆变换分布}可知
\begin{align}
q(y_1,\cdots,y_n)=
\begin{cases}
p(x_1(y_1,\cdots,y_n),\cdots,x_n(y_1,\cdots,y_n))\,|J|, & \text{若 }(y_1,\cdots,y_n)\text{ 属于 }g_1,\cdots,g_n\text{ 的值域},\\
0, & \text{其他}.
\end{cases}\label{align:ch03-变换密度公式}
\end{align}
其中 $J$ 为坐标变换的 Jacobi 行列式
\begin{align}
J=\left|\begin{array}{ccc}
\dfrac{\partial x_1}{\partial y_1} & \cdots & \dfrac{\partial x_1}{\partial y_n}\\
\vdots & & \vdots\\
\dfrac{\partial x_n}{\partial y_1} & \cdots & \dfrac{\partial x_n}{\partial y_n}
\end{array}\right|.\label{align:ch03-Jacobi}
\end{align}
这里，我们假定上述偏导数存在而且连续。

公式\eqref{align:ch03-变换密度公式}对应于单变量场合的公式\eqref{align:ch03-单调变换密度}，也可以导出对应于公式\eqref{align:ch03-分段单调密度}的多变量场合的公式，这留给读者作为练习。

\begin{example}[线性变换]\label{example:ch03-二维线性变换}
若 $(\xi_1,\xi_2)$ 的密度函数为 $p(x_1,x_2)$，而
\begin{align}
\begin{aligned}
\eta_1&=a\xi_1+b\xi_2,\\
\eta_2&=c\xi_1+d\xi_2,
\end{aligned}\label{align:ch03-二维线性变换}
\end{align}
这里 $\Delta=\begin{vmatrix}a&b\\ c&d\end{vmatrix}\neq 0$，试求 $(\eta_1,\eta_2)$ 的密度函数 $q(y_1,y_2)$。

\paragraph*{解}在本例中
\begin{align*}y_1=g_1(x_1,x_2)=ax_1+bx_2,\qquad y_2=g_2(x_1,x_2)=cx_1+dx_2.\end{align*}
因此
\begin{align*}x_1=\frac{d}{\Delta}y_1-\frac{b}{\Delta}y_2,\qquad x_2=-\frac{c}{\Delta}y_1+\frac{a}{\Delta}y_2.\end{align*}
而
\begin{align*}J=\left|\begin{array}{rr}
\dfrac{d}{\Delta} & -\dfrac{b}{\Delta}\\[4pt]
-\dfrac{c}{\Delta} & \dfrac{a}{\Delta}
\end{array}\right|=\frac{ad-bc}{\Delta^{2}}=\frac{1}{ad-bc}.\end{align*}
最后得到
\begin{align}
q(y_1,y_2)=\frac{p\!\left(\dfrac{d}{\Delta}y_1-\dfrac{b}{\Delta}y_2,\,-\dfrac{c}{\Delta}y_1+\dfrac{a}{\Delta}y_2\right)}{|ad-bc|}.\label{align:ch03-二维线性变换密度}
\end{align}
\end{example}

\begin{example}[$\chi^{2}$ 分布的可加性与 $F$ 分布]\label{example:ch03-卡方可加}
若 $\xi$ 与 $\eta$ 相互独立，分别服从自由度为 $m$ 和 $n$ 的 $\chi^{2}$ 分布，试求 $\alpha=\xi+\eta$ 与 $\beta=\dfrac{\xi}{\eta}\cdot\dfrac{n}{m}$ 的密度函数 $q(u,v)$。

\paragraph*{解}由题设据\eqref{align:ch03-卡方密度}知 $(\xi,\eta)$ 的联合密度函数为
\begin{align*}p(x,y)=\frac{1}{2^{\frac{m+n}{2}}\Gamma\!\left(\frac{m}{2}\right)\Gamma\!\left(\frac{n}{2}\right)}x^{\frac{m}{2}-1}y^{\frac{n}{2}-1}\mathrm{e}^{-\frac{x+y}{2}},\qquad x>0,\ y>0.\end{align*}

对 $u>0,v>0$ 作变换 $u=x+y$，$v=\dfrac{x}{y}\cdot\dfrac{n}{m}$，其逆变换为
\begin{align*}x=\frac{muv}{n+mv},\qquad y=\frac{nu}{n+mv}.\end{align*}
由于
\begin{align*}
J^{-1}=\left|\begin{array}{cc}
\dfrac{\partial u}{\partial x} & \dfrac{\partial u}{\partial y}\\[4pt]
\dfrac{\partial v}{\partial x} & \dfrac{\partial v}{\partial y}
\end{array}\right|=\left|\begin{array}{cc}
1 & 1\\
\dfrac{n}{ym} & -\dfrac{xn}{y^{2}m}
\end{array}\right|=-\frac{n(x+y)}{my^{2}}
=-\frac{n}{m}\cdot\frac{\left(1+\dfrac{m}{n}v\right)^{2}}{u}.
\end{align*}
因此
\begin{align*}|J|=\frac{m}{n}\cdot\frac{u}{\left(1+\dfrac{m}{n}v\right)^{2}}.\end{align*}
于是 $(\alpha,\beta)$ 的联合分布密度函数为
\begin{align*}
q(u,v)&=\frac{1}{2^{\frac{m+n}{2}}\Gamma\!\left(\frac{m}{2}\right)\Gamma\!\left(\frac{n}{2}\right)}\mathrm{e}^{-\frac{u}{2}}\left(\frac{m}{n}\right)^{\frac{m}{2}-1}u^{\frac{m+n}{2}-2}\\
&\quad\times\frac{v^{\frac{m}{2}-1}}{\left(1+\dfrac{m}{n}v\right)^{\frac{m+n}{2}-2}}\cdot\frac{m}{n}\cdot\frac{u}{\left(1+\dfrac{m}{n}v\right)^{2}}\\
&=\frac{1}{2^{\frac{m+n}{2}}\Gamma\!\left(\frac{m+n}{2}\right)}u^{\frac{m+n}{2}-1}\mathrm{e}^{-\frac{u}{2}}\\
&\quad\times\frac{\Gamma\!\left(\frac{m+n}{2}\right)\cdot\left(\frac{m}{n}\right)^{\frac{m}{2}}v^{\frac{m}{2}-1}}{\Gamma\!\left(\frac{m}{2}\right)\Gamma\!\left(\frac{n}{2}\right)\cdot\left(1+\dfrac{m}{n}v\right)^{\frac{m+n}{2}}}.
\end{align*}

因此，$\alpha$ 与 $\beta$ 独立，而且 $\alpha$ 服从自由度为 $m+n$ 的 $\chi^{2}$ 分布。这揭示了 $\chi^{2}$ 分布的一个重要性质：若两个相互独立的 $\chi^{2}$ 分布随机变量 $\xi$ 与 $\eta$，它们各具有自由度 $m$ 及 $n$，则其和 $\xi+\eta$ 服从自由度为 $m+n$ 的 $\chi^{2}$ 分布，这个性质称为 $\chi^{2}$ 分布的\textbf{可加性}，简记作
\begin{align}
\chi_m^{2}*\chi_n^{2}=\chi_{m+n}^{2}.\label{align:ch03-卡方可加}
\end{align}

利用例\ref{example:ch03-平方变换}的结果和这里证明了的 $\chi^{2}$ 分布的可加性，不难给出 $\chi^{2}$ 分布的一种推导。事实上，$\chi_n^{2}$ 分布是作为 $n$ 个相互独立的标准正态变量的平方和的分布而命名的，明白这点之后，会觉得关于它的那些结论都是相当自然的。$\chi^{2}$ 分布的直接推导见习题 $37$。

随机变量 $\beta=\dfrac{\xi}{\eta}\cdot\dfrac{n}{m}$ 的密度函数为
\begin{align}
f(x;m,n)=
\begin{cases}
\dfrac{\Gamma\!\left(\frac{m+n}{2}\right)}{\Gamma\!\left(\frac{m}{2}\right)\Gamma\!\left(\frac{n}{2}\right)}\left(\dfrac{m}{n}\right)^{\frac{m}{2}}x^{\frac{m}{2}-1}\\[4pt]
\quad\div\left(1+\dfrac{m}{n}x\right)^{\frac{m+n}{2}}, & x>0,\\[4pt]
0, & x\leqslant 0.
\end{cases}\label{align:ch03-F分布}
\end{align}
这个分布称为 $F$ 分布，是数理统计中重要分布之一。
\end{example}

在讨论随机向量的变换中，着重研究有一一对应变换的情况是很自然的。如果 $m>n$，则当 $(\xi_1,\cdots,\xi_n)$ 有密度函数时，$(\eta_1,\cdots,\eta_m)$ 不会有密度函数存在，不能深入讨论。但是当 $m<n$ 时，则有可能利用变换法，关键是增补变量，使之成为一一对应的情况。

在上一小节讨论随机向量的函数的分布律中，主要只讲直接法，下面我们通过一个例子来说明在这种场合如何通过增补变量使用变换法。

\begin{example}[$t$ 分布]\label{example:ch03-t分布}
设 $\xi,\eta$ 为两个独立随机变量，$\xi$ 服从 $N(0,1)$，$\eta$ 服从自由度为 $n$ 的 $\chi^{2}$ 分布\eqref{align:ch03-卡方密度}，令 $T=\xi\bigg/\sqrt{\dfrac{\eta}{n}}$，试求 $T$ 的密度函数。

\paragraph*{解}为求得 $T$ 的密度函数，引进增补变量 $S=\eta$，先求 $(S,T)$ 的联合密度函数。

$\xi,\eta$ 相互独立，故 $(\xi,\eta)$ 的联合密度函数为
\begin{align*}p(x,y)=\frac{1}{\sqrt{2\pi}}\mathrm{e}^{-\frac{x^{2}}{2}}\cdot\frac{1}{2^{n/2}\Gamma\!\left(\frac{n}{2}\right)}y^{\frac{n}{2}-1}\mathrm{e}^{-\frac{y}{2}},\qquad -\infty<x<\infty,\ y>0.\end{align*}
变换 $s=y$，$t=\dfrac{x}{\sqrt{y/n}}$ 的逆变换为 $x=t\left(\dfrac{s}{n}\right)^{1/2}$，$y=s$，其雅可比行列式
\begin{align*}J=\begin{vmatrix}
\dfrac{\partial x}{\partial s} & \dfrac{\partial x}{\partial t}\\[4pt]
\dfrac{\partial y}{\partial s} & \dfrac{\partial y}{\partial t}
\end{vmatrix}=\begin{vmatrix}
\dfrac{t}{2n}\left(\dfrac{s}{n}\right)^{-\frac{1}{2}} & \left(\dfrac{s}{n}\right)^{\frac{1}{2}}\\
1 & 0
\end{vmatrix}=-\left(\dfrac{s}{n}\right)^{\frac{1}{2}}.\end{align*}
\begin{align*}|J|=\left(\dfrac{s}{n}\right)^{\frac{1}{2}}.\end{align*}
故 $(S,T)$ 的联合密度函数为
\begin{align*}
q(s,t)&=p\left(t\left(\frac{s}{n}\right)^{\frac{1}{2}},s\right)|J|\\
&=\frac{1}{\sqrt{2\pi}}\mathrm{e}^{-\frac{s t^{2}}{2n}}\cdot\frac{1}{2^{\frac{n}{2}}\Gamma\!\left(\frac{n}{2}\right)}s^{\frac{n}{2}-1}\mathrm{e}^{-\frac{s}{2}}\cdot\left(\frac{s}{n}\right)^{\frac{1}{2}}.
\end{align*}
因而 $T$ 的密度函数为
\begin{align*}
p_T(t)=\int_{0}^{\infty}q(s,t)\,\mathrm{d}s
&=\frac{1}{2^{\frac{n+1}{2}}\sqrt{n\pi}\,\Gamma\!\left(\frac{n}{2}\right)}\int_{0}^{\infty}\mathrm{e}^{-\left(1+\frac{t^{2}}{n}\right)\frac{s}{2}}s^{\frac{n+1}{2}-1}\,\mathrm{d}s\\
&=\frac{\left(1+\dfrac{t^{2}}{n}\right)^{-\frac{n+1}{2}}}{\sqrt{n\pi}\,\Gamma\!\left(\frac{n}{2}\right)}\int_{0}^{\infty}\mathrm{e}^{-u}u^{\frac{n+1}{2}-1}\,\mathrm{d}u\\
&=\frac{\Gamma\!\left(\frac{n+1}{2}\right)}{\sqrt{n\pi}\,\Gamma\!\left(\frac{n}{2}\right)}\left(1+\frac{t^{2}}{n}\right)^{-\frac{n+1}{2}}.
\end{align*}
密度函数
\begin{align}
t(x;n)=\frac{\Gamma\!\left(\frac{n+1}{2}\right)}{\sqrt{n\pi}\,\Gamma\!\left(\frac{n}{2}\right)}\left(\frac{x^{2}}{n}+1\right)^{-(n+1)/2}\label{align:ch03-t分布密度}
\end{align}
称为自由度为 $n$ 的 $t$ 分布，它是数理统计中另一重要分布。
\end{example}

\subsection*{五、随机变量的函数的独立性}

首先证明一个定理。

\begin{theorem}[独立随机变量的函数仍独立]\label{theorem:ch03-独立函数仍独立}
若 $\xi_1,\cdots,\xi_n$ 是相互独立的随机变量，则 $f_1(\xi_1),\cdots,f_n(\xi_n)$ 也是相互独立的，这里 $f_i$（$i=1,\cdots,n$）是任意的一元 Borel 函数。
\end{theorem}

\begin{proof}
对任意的一维 Borel 点集 $A_1,\cdots,A_n$ 有
\begin{align*}
P\{f_1(\xi_1)\in A_1,\cdots,f_n(\xi_n)\in A_n\}
&=P\{\xi_1\in f_1^{-1}(A_1),\cdots,\xi_n\in f_n^{-1}(A_n)\}\\
&=P\{\xi_1\in f_1^{-1}(A_1)\}\cdots P\{\xi_n\in f_n^{-1}(A_n)\}\\
&=P\{f_1(\xi_1)\in A_1\}\cdots P\{f_n(\xi_n)\in A_n\}.
\end{align*}
\end{proof}

定理的结论在直观上是明显的，但在定理的证明中却要两次用到未证明的论断\eqref{align:ch03-独立Borel}，其中第一次用来指明对 $\xi_1,\cdots,\xi_n$ 的有关概率可以化为乘积的形式，另一次用来说明最后的等式表明 $f_1(\xi_1),\cdots,f_n(\xi_n)$ 是相互独立的，且第一次是难以避免的。

这个结果可以推广到随机向量的场合。

例\ref{example:ch03-卡方可加}说明，即使由相同的随机向量构成的不同函数也可能是独立的，这种情况在概率论与数理统计中相当重要，下面再讨论一些例子。

\begin{example}[极坐标变换]\label{example:ch03-极坐标变换}
若 $\xi$ 与 $\eta$ 是相互独立的随机变量，均服从 $N(0,1)$，试证化为极坐标后，$\rho=\sqrt{\xi^{2}+\eta^{2}}$ 与 $\varphi=\operatorname{arctg}\!\left(\dfrac{\eta}{\xi}\right)$（$\varphi$ 取值于 $[0,2\pi]$），是相互独立的。

\paragraph*{解}采用极坐标，$x=r\cos\theta$，$y=r\sin\theta$，因此 $r=\sqrt{x^{2}+y^{2}}$，$\theta=\operatorname{arctg}\dfrac{y}{x}$，因为 $(\xi,\eta)$ 的密度函数为
\begin{align*}p(x,y)=\frac{1}{2\pi}\mathrm{e}^{-(x^{2}+y^{2})/2},\end{align*}
而
\begin{align*}J=\begin{vmatrix}
\dfrac{\partial x}{\partial r} & \dfrac{\partial x}{\partial\theta}\\[4pt]
\dfrac{\partial y}{\partial r} & \dfrac{\partial y}{\partial\theta}
\end{vmatrix}=\begin{vmatrix}
\cos\theta & -r\sin\theta\\
\sin\theta & r\cos\theta
\end{vmatrix}=r.\end{align*}
故 $(\rho,\varphi)$ 的密度函数为
\begin{align*}
q(r,\theta)&=\frac{1}{2\pi}\mathrm{e}^{-(x^{2}+y^{2})/2}\cdot r\\
&=\frac{1}{2\pi}\cdot r\mathrm{e}^{-r^{2}/2},\qquad r\geqslant 0,\ 0\leqslant\theta\leqslant 2\pi.
\end{align*}
即 $\rho=\sqrt{\xi^{2}+\eta^{2}}$ 的密度函数为
\begin{align}
R(r)=
\begin{cases}
r\mathrm{e}^{-r^{2}/2}, & r\geqslant 0,\\
0, & r<0.
\end{cases}\label{align:ch03-rayleigh密度}
\end{align}
这个分布称为 Rayleigh 分布。

而 $\varphi=\operatorname{arctg}\dfrac{\eta}{\xi}$ 服从 $[0,2\pi]$ 中的均匀分布，并且 $\rho$ 与 $\varphi$ 是独立的。

这个结果常被用来产生服从正态分布的随机数。做法如下：产生相互独立的 $[0,1]$ 均匀分布的随机数 $U_1,U_2$，令
\begin{align}
\begin{aligned}
\xi&=(-2\ln U_1)^{\frac{1}{2}}\cos 2\pi U_2,\\
\eta&=(-2\ln U_1)^{\frac{1}{2}}\sin 2\pi U_2.
\end{aligned}\label{align:ch03-boxmuller}
\end{align}
则 $\xi$ 与 $\eta$ 是相互独立的 $N(0,1)$ 随机数。

这样做法的理由让读者自行论证。
\end{example}

\begin{example}[二维正态的旋转变换]\label{example:ch03-正态旋转}
若 $(\xi_1,\xi_2)$ 服从二元正态分布\eqref{align:ch03-二元正态密度}，其中 $\mu_1=\mu_2=0$。令
\begin{align*}\eta_1=\xi_1\cos\alpha+\xi_2\sin\alpha,\qquad \eta_2=-\xi_1\sin\alpha+\xi_2\cos\alpha.\end{align*}
这里 $0\leqslant\alpha\leqslant 2\pi$，是某个角度。我们来求 $(\eta_1,\eta_2)$ 的密度函数 $q(u,v)$。

这里可直接用例\ref{example:ch03-二维线性变换}的结果。其中 $J=1$，因此
\begin{align}
q(u,v)=p(u\cos\alpha-v\sin\alpha,u\sin\alpha+v\cos\alpha)
=\frac{1}{2\pi\sigma_1\sigma_2\sqrt{1-\rho^{2}}}\exp\left\{-\frac{1}{2(1-\rho^{2})}(Au^{2}-2Buv+Cv^{2})\right\},\label{align:ch03-旋转密度}
\end{align}
其中
\begin{align*}
A&=\frac{\cos^{2}\alpha}{\sigma_1^{2}}-2\rho\frac{\cos\alpha\sin\alpha}{\sigma_1\sigma_2}+\frac{\sin^{2}\alpha}{\sigma_2^{2}},\\
B&=\frac{\cos\alpha\sin\alpha}{\sigma_1^{2}}-\rho\frac{\sin^{2}\alpha-\cos^{2}\alpha}{\sigma_1\sigma_2}-\frac{\cos\alpha\sin\alpha}{\sigma_2^{2}},\\
C&=\frac{\sin^{2}\alpha}{\sigma_1^{2}}+2\rho\frac{\cos\alpha\sin\alpha}{\sigma_1\sigma_2}+\frac{\cos^{2}\alpha}{\sigma_2^{2}}.
\end{align*}
由\eqref{align:ch03-旋转密度}可看出由二维正态向量 $(\xi_1,\xi_2)$ 经坐标旋转而得的随机向量 $(\eta_1,\eta_2)$ 还是服从正态分布。进一步，若选 $\alpha$ 使得
\begin{align}
\operatorname{tg}2\alpha=\frac{2\rho\sigma_1\sigma_2}{\sigma_1^{2}-\sigma_2^{2}},\label{align:ch03-旋转角}
\end{align}
则 $B=0$，因此 $\eta_1$ 与 $\eta_2$ 独立。这说明二元正态分布密度可经适当的坐标旋转化为两个正态分布密度之积。利用正交变换把多维正态变量化作独立正态分量，在数理统计中有重要应用。
\end{example}

\end{document}